\documentclass[11pt,a4paper]{article}

\usepackage[utf8]{inputenc}
\usepackage[T1]{fontenc}
\usepackage{lmodern}
\usepackage[a4paper,margin=1in]{geometry}
\usepackage{amsmath,amssymb,mathtools}
\usepackage{booktabs}
\usepackage{enumitem}
\usepackage[hidelinks]{hyperref}

\title{A Candidate Joint Hierarchical Bayesian Logistic Regression Model}
\author{}
\date{\today}

\begin{document}
\maketitle

\section{Purpose and notation}

The aim is to estimate the risk of death or diagnosis and to quantify its uncertainty using a Bayesian model. We compare a two-stage approximation with a joint model that represents uncertainty in both the old-data and new-data components.
1
Let the datasets be
\[
D_{\mathrm{old}}=\{(y_i^{(o)},x_i^{(o)})\}_{i=1}^{n_o},
\qquad
D_{\mathrm{new}}=\{(y_j^{(n)},x_j^{(n)})\}_{j=1}^{n_n},
\]
where $y\in\{0,1\}$ is the binary outcome and $x$ is a vector of predictors. Write $\sigma(z)=1/(1+\exp(-z))$ for the logistic function.

\textbf{Important:} The specification below is a candidate template. The correct joint model depends on the likelihood and parameter relationships in the implemented two-stage model. In particular, if the datasets contain overlapping/repeated individuals, their dependence must be modelled rather than assuming independent likelihoods.

\section{Likelihood for the old dataset}

For observations in $D_{\mathrm{old}}$, suppose
\[
Y_i^{(o)}\mid \beta \sim \operatorname{Bernoulli}(p_i^{(o)}),
\qquad
p_i^{(o)}=\sigma\!\left((x_i^{(o)})^\top\beta\right).
\]
The old-data likelihood is
\begin{align}
L_o(\beta)
&=p(D_{\mathrm{old}}\mid\beta)\\
&=\prod_{i=1}^{n_o}
\left(p_i^{(o)}\right)^{y_i^{(o)}}
\left(1-p_i^{(o)}\right)^{1-y_i^{(o)}}.
\end{align}

\section{Likelihood for the new dataset}

For observations in $D_{\mathrm{new}}$, suppose
\[
Y_j^{(n)}\mid \theta \sim \operatorname{Bernoulli}(p_j^{(n)}),
\qquad
p_j^{(n)}=\sigma\!\left((x_j^{(n)})^\top\theta\right).
\]
The new-data likelihood is
\begin{align}
L_n(\theta)
&=p(D_{\mathrm{new}}\mid\theta)\\
&=\prod_{j=1}^{n_n}
\left(p_j^{(n)}\right)^{y_j^{(n)}}
\left(1-p_j^{(n)}\right)^{1-y_j^{(n)}}.
\end{align}

\section{Priors and hierarchical relationship}

One possible hierarchical specification is
\[
\beta\sim\mathcal N(m_0,\Sigma_0),
\]
and
\[
\theta\mid\beta,\tau\sim\mathcal N(A\beta,\tau^2 I),
\qquad \tau>0.
\]
If $\tau$ is unknown, assign it a proper prior $p(\tau)$. Here $A$ is a specified matrix that maps the old-model coefficients to the scale and dimension of the new-model coefficients. Its use must be justified for the application; if the coefficients do not have a meaningful mapping, a different hierarchical link is needed.

Conditional on the parameters, and assuming the datasets are independent,
\[
p(D_{\mathrm{old}},D_{\mathrm{new}}\mid\beta,\theta)
=L_o(\beta)L_n(\theta).
\]

\section{Joint posterior derivation}

By Bayes' rule,
\[
p(\beta,\theta,\tau\mid D_{\mathrm{old}},D_{\mathrm{new}})
\propto
p(D_{\mathrm{old}},D_{\mathrm{new}}\mid\beta,\theta,\tau)
\,p(\theta\mid\beta,\tau)\,p(\beta)\,p(\tau).
\]
Using the conditional-independence assumption and the priors above gives
\begin{align}
&p(\beta,\theta,\tau\mid D_{\mathrm{old}},D_{\mathrm{new}})\\
&\quad\propto
L_o(\beta)L_n(\theta)\,
\mathcal N(\theta;A\beta,\tau^2 I)\,
\mathcal N(\beta;m_0,\Sigma_0)\,
p(\tau).
\label{eq:joint-posterior}
\end{align}
The proportionality omits the marginal likelihood, which is a normalising constant with respect to $(\beta,\theta,\tau)$.

If $\tau$ is fixed rather than unknown, remove $\tau$ from the parameter set and omit $p(\tau)$ from \eqref{eq:joint-posterior}.

\subsection{Log-posterior up to an additive constant}

For computation, the log-posterior can be written as
\begin{align}
\log p(\beta,\theta,\tau\mid D_{\mathrm{old}},D_{\mathrm{new}})
={}&
\sum_{i=1}^{n_o}\left[
y_i^{(o)}\log p_i^{(o)}
+(1-y_i^{(o)})\log(1-p_i^{(o)})
\right]\\
&+\sum_{j=1}^{n_n}\left[
y_j^{(n)}\log p_j^{(n)}
+(1-y_j^{(n)})\log(1-p_j^{(n)})
\right]\\
&+\log\mathcal N(\theta;A\beta,\tau^2I)
+\log\mathcal N(\beta;m_0,\Sigma_0)
+\log p(\tau)+C,
\end{align}
where $C$ does not depend on the unknown parameters. Include or omit normal-distribution constants consistently depending on whether the expression is used only up to proportionality or for a fully normalised density.

\section{Risk estimation}

For a target covariate vector $x_\star$, define the risk under the new-data model as
\[
\omega_\star=\Pr(Y_\star=1\mid x_\star,\theta)
=\sigma(x_\star^\top\theta).
\]
The posterior distribution of the risk is induced by the joint posterior:
\begin{align}
p(\omega_\star\mid D_{\mathrm{old}},D_{\mathrm{new}},x_\star)
=\int &\delta\!\left(\omega_\star-\sigma(x_\star^\top\theta)\right)\\
&\times p(\beta,\theta,\tau\mid D_{\mathrm{old}},D_{\mathrm{new}})
\,d\beta\,d\theta\,d\tau,
\end{align}
where $\delta(\cdot)$ is the Dirac delta distribution. Equivalently, draw $(\beta^{(s)},\theta^{(s)},\tau^{(s)})$ from the joint posterior and calculate
\[
\omega_\star^{(s)}=\sigma(x_\star^\top\theta^{(s)}).
\]
The posterior draws $\{\omega_\star^{(s)}\}_{s=1}^S$ provide a posterior mean or median and credible intervals for the risk.

\section{Comparison with the two-stage approximation}

A two-stage approach may fit the old model first, then use a point estimate such as $\widehat\beta$ as fixed input to the new-stage model. This can omit uncertainty in the old-stage estimate. The joint model instead integrates over uncertainty in $\beta$ and any hyperparameters, conditional on the specified model.

To compare the approaches fairly, evaluate both on the same held-out observations or cross-validation splits. Possible metrics include calibration, Brier score, predictive log score, and the coverage and width of predictive or credible intervals. The comparison should distinguish predictive accuracy from uncertainty calibration.

\section{Assumptions to verify before using this model}

\begin{enumerate}[label=\arabic*.]
  \item The outcomes and predictors are correctly defined for both datasets.
  \item The two likelihoods are appropriate and the datasets can be treated as conditionally independent, or their dependence is explicitly modelled.
  \item The relationship $\theta\mid\beta,\tau\sim\mathcal N(A\beta,\tau^2I)$ is scientifically justified, including the choice of $A$ and the scale parameter.
  \item The priors are proper and suitable for the parameter scales.
  \item Any systematic bias in the old model is represented explicitly if the scientific setup calls for it; a hierarchical prior does not automatically remove bias.
  \item The target risk $\omega_\star$ corresponds to the intended population, outcome, and prediction time horizon.
\end{enumerate}

\end{document}
